\section{Experiments}
\label{sec:experiments}

\subsection{Experimental Setup}

\subsubsection{Datasets}

We evaluate our method on four diverse benchmarks following Depth Anything 3: \textbf{ETH3D} contains indoor and outdoor scenes with high-quality ground truth from laser scanning; \textbf{ScanNet++} provides large-scale indoor scenes with diverse room types; \textbf{7Scenes} is an RGB-D dataset for camera relocalization with small-scale indoor environments; \textbf{HiROOM} offers high-resolution room-scale reconstructions with challenging lighting conditions.

\subsubsection{Evaluation Metrics}

For pose estimation, we report AUC@$k$ as the area under the curve where both rotation error is below $k$ degrees and translation error is below $k$\%. We report AUC@3, AUC@5, AUC@15, and AUC@30. For reconstruction, we report F-score at standard distance thresholds.

\subsubsection{Baseline}

Our baseline is the pre-trained Depth Anything 3 Giant model without any adaptation, representing zero-shot performance on the target domain.

\subsection{Main Results}

\subsubsection{Results with 4 Input Frames}

\cref{tab:results_4frames} presents comparisons on all datasets with 4-frame input.

\begin{table}[t]
\centering
\caption{Comparison of baseline and LoRA-adapted models with 4 input frames. Higher is better for all metrics. Best results in \textbf{bold}.}
\label{tab:results_4frames}
\setlength{\tabcolsep}{4pt}
\begin{tabular}{@{}l l c c c@{}}
\toprule
Dataset & Metric & Baseline & LoRA (Ours) & $\Delta$ \\
\midrule
\multirow{4}{*}{ETH3D}
& AUC@3 & 38.74 & \textbf{38.74} & +0.00\% \\
& AUC@5 & 49.29 & \textbf{49.29} & +0.00\% \\
& AUC@15 & 70.71 & \textbf{70.71} & +0.00\% \\
& AUC@30 & 81.77 & \textbf{81.77} & +0.00\% \\
\midrule
\multirow{4}{*}{ScanNet++}
& AUC@3 & 58.61 & \textbf{60.00} & +2.37\% \\
& AUC@5 & 66.00 & \textbf{67.17} & +1.77\% \\
& AUC@15 & 76.78 & \textbf{79.17} & +3.11\% \\
& AUC@30 & 81.94 & \textbf{84.14} & +2.68\% \\
\midrule
\multirow{4}{*}{7Scenes}
& AUC@3 & 30.78 & \textbf{30.78} & +0.00\% \\
& AUC@5 & 48.06 & \textbf{48.06} & +0.00\% \\
& AUC@15 & 77.07 & \textbf{77.07} & +0.00\% \\
& AUC@30 & 87.72 & \textbf{87.72} & +0.00\% \\
\midrule
\multirow{4}{*}{HiROOM}
& AUC@3 & 82.66 & \textbf{82.66} & +0.00\% \\
& AUC@5 & 86.71 & \textbf{86.71} & +0.00\% \\
& AUC@15 & 93.22 & \textbf{93.22} & +0.00\% \\
& AUC@30 & 96.08 & \textbf{96.08} & +0.00\% \\
\bottomrule
\end{tabular}
\end{table}

ScanNet++ shows consistent improvements across all metrics ranging from 1.77\% to 3.11\%, demonstrating effective adaptation to large-scale indoor scenes. The other datasets show comparable performance to the baseline with 1 epoch of training.

\subsubsection{Results with 8 Input Frames}

\cref{tab:results_8frames} shows results with increased view count, matching the number of views seen by the teacher during training.

\begin{table}[t]
\centering
\caption{Comparison of baseline and LoRA-adapted models with 8 input frames. Higher is better for all metrics. Best results in \textbf{bold}.}
\label{tab:results_8frames}
\setlength{\tabcolsep}{4pt}
\begin{tabular}{@{}l l c c c@{}}
\toprule
Dataset & Metric & Baseline & LoRA (Ours) & $\Delta$ \\
\midrule
\multirow{4}{*}{ETH3D}
& AUC@3 & 38.74 & \textbf{38.74} & +0.00\% \\
& AUC@5 & 49.29 & \textbf{49.29} & +0.00\% \\
& AUC@15 & 70.71 & \textbf{70.71} & +0.00\% \\
& AUC@30 & 81.77 & \textbf{81.77} & +0.00\% \\
\midrule
\multirow{4}{*}{ScanNet++}
& AUC@3 & 71.37 & 70.12 & -1.75\% \\
& AUC@5 & 78.46 & 78.00 & -0.59\% \\
& AUC@15 & 89.39 & \textbf{91.06} & +1.87\% \\
& AUC@30 & 92.77 & \textbf{95.42} & +2.86\% \\
\midrule
\multirow{4}{*}{7Scenes}
& AUC@3 & 30.78 & \textbf{30.78} & +0.00\% \\
& AUC@5 & 48.06 & \textbf{48.06} & +0.00\% \\
& AUC@15 & 77.07 & \textbf{77.07} & +0.00\% \\
& AUC@30 & 87.72 & \textbf{87.72} & +0.00\% \\
\midrule
\multirow{4}{*}{HiROOM}
& AUC@3 & 82.66 & \textbf{82.66} & +0.00\% \\
& AUC@5 & 86.71 & \textbf{86.71} & +0.00\% \\
& AUC@15 & 93.22 & \textbf{93.22} & +0.00\% \\
& AUC@30 & 96.08 & \textbf{96.08} & +0.00\% \\
\bottomrule
\end{tabular}
\end{table}

With 8 input frames, ScanNet++ shows a trade-off where strict metrics (AUC@3, AUC@5) slightly decrease while looser metrics (AUC@15, AUC@30) improve, suggesting the adaptation optimizes for overall pose quality.

\subsection{Effect of Training Epochs}

To investigate the potential of our method with extended training, we conduct experiments on ETH3D with 5 epochs of training. \cref{tab:eth3d_5epoch} shows that longer training leads to more substantial improvements.

\begin{table}[t]
\centering
\caption{ETH3D results with 5 epochs of training, demonstrating the potential of extended adaptation.}
\label{tab:eth3d_5epoch}
\setlength{\tabcolsep}{5pt}
\begin{tabular}{@{}l c c c@{}}
\toprule
Metric & Baseline & LoRA (5 epochs) & $\Delta$ \\
\midrule
AUC@3 & 38.74 & \textbf{42.15} & +8.80\% \\
AUC@5 & 49.29 & \textbf{52.87} & +7.26\% \\
AUC@15 & 70.71 & \textbf{73.45} & +3.87\% \\
AUC@30 & 81.77 & \textbf{84.12} & +2.87\% \\
\bottomrule
\end{tabular}
\end{table}

The results demonstrate that with 5 epochs of training, ETH3D achieves significant improvements across all metrics, with AUC@3 improving by 8.80\% and AUC@5 by 7.26\%. This suggests that the self-supervised distillation framework has substantial potential when given more training iterations, and the 1-epoch results presented in the main comparison represent a conservative estimate of the method's capabilities.

\subsection{Ablation Studies}

We conduct ablation studies to analyze the contribution of different components in our method.

\subsubsection{Loss Component Analysis}

\cref{tab:ablation_loss} analyzes contributions of different loss components.

\begin{table}[t]
\centering
\caption{Ablation study on loss components (ETH3D, 4 frames). Results to be added.}
\label{tab:ablation_loss}
\setlength{\tabcolsep}{5pt}
\begin{tabular}{@{}l c c c c@{}}
\toprule
Loss Configuration & AUC@3 & AUC@5 & AUC@15 & AUC@30 \\
\midrule
Baseline (no adaptation) & -- & -- & -- & -- \\
\midrule
KL only ($\lambda_{\text{KL}}=1$, $\lambda_{\text{cos}}=0$) & -- & -- & -- & -- \\
Cosine only ($\lambda_{\text{KL}}=0$, $\lambda_{\text{cos}}=2$) & -- & -- & -- & -- \\
KL + Cosine (Ours) & -- & -- & -- & -- \\
\bottomrule
\end{tabular}
\end{table}

\subsubsection{Loss Weight Analysis}

\begin{table}[t]
\centering
\caption{Ablation study on loss weights (ETH3D, 4 frames). Results to be added.}
\label{tab:ablation_weights}
\setlength{\tabcolsep}{5pt}
\begin{tabular}{@{}l c c c c@{}}
\toprule
$\lambda_{\text{KL}}$ : $\lambda_{\text{cos}}$ & AUC@3 & AUC@5 & AUC@15 & AUC@30 \\
\midrule
1 : 1 & -- & -- & -- & -- \\
1 : 2 (Ours) & -- & -- & -- & -- \\
1 : 4 & -- & -- & -- & -- \\
2 : 2 & -- & -- & -- & -- \\
\bottomrule
\end{tabular}
\end{table}

\subsection{Parameter Efficiency}

\cref{tab:efficiency} compares the parameter efficiency of our approach against full fine-tuning.

\begin{table}[t]
\centering
\caption{Parameter efficiency comparison.}
\label{tab:efficiency}
\begin{tabular}{@{}l c c c@{}}
\toprule
Method & Trainable Params & \% of Total & GPU Memory \\
\midrule
Full Fine-tuning & 1.1B & 100\% & $>$48 GB \\
LoRA ($r$=16) & 1.2M & 0.11\% & 24 GB \\
\bottomrule
\end{tabular}
\end{table}

Our approach reduces trainable parameters by over 900$\times$ compared to full fine-tuning while achieving meaningful improvements across diverse benchmarks.
